\subsection{Microservicio de Identidad y Académico}
\label{subsec:identidad-academico}

\subsubsection{Alcance y responsabilidad}

El microservicio controla quién accede a CLAB y con qué permisos, y mantiene
la información académica con la que se vinculan las reservas. Su motivo de
cambio son las reglas de acceso y la organización académica de la
institución. El alcance se definió a partir del código del sistema y de las
historias de usuario de la Entrega~1, analizando el dominio y sus
dependencias para identificar los límites del servicio
\cite{abgaz2023decomposition,oumoussa2024evolution,velepucha2023survey}.

El servicio atiende el inicio de sesión con credenciales institucionales
(RF-01), la asignación de roles y permisos a las cuentas (RF-02) y la gestión
de carreras, asignaturas, horarios y periodos académicos (RF-16). También
administra las solicitudes de acceso de nuevos usuarios y emite el token con
el que el API Gateway identifica a cada usuario. Para el RF-14, registra sus
propios cambios y los publica; Información y Avisos reúne la vista central.

\subsubsection{Propiedad de datos}

El servicio usa una base PostgreSQL propia, con las siguientes entidades:

\begin{itemize}
    \item \textbf{Usuario y sesión activa.} El usuario conserva sus
    credenciales cifradas, su estado y sus datos de contacto. La sesión
    activa registra el token emitido y permite cerrarlo al finalizar la
    sesión.
    \item \textbf{Rol, permiso y módulo.} Los roles agrupan permisos sobre
    los módulos del sistema. Las relaciones entre usuario y rol, y entre rol
    y permiso, determinan qué operaciones puede ejecutar cada cuenta.
    \item \textbf{Solicitud de acceso.} Conserva la petición de un nuevo
    usuario, su estado (pendiente, aprobada o rechazada) y quién la resolvió.
    \item \textbf{Carrera, asignatura, horario y periodo académico.}
    Las asignaturas pertenecen a una carrera; los horarios vinculan una
    asignatura con un docente y un periodo.
    \item \textbf{Estudiante, matrícula e inscripción.} Relacionan a los
    estudiantes con las carreras y asignaturas del periodo vigente.
    \item \textbf{Sede y facultad.} Mantienen la organización institucional
    que Laboratorios y Recursos referencia por identificador.
    % Confirmar con el grupo que Sede y Facultad quedan en este servicio.
    \item \textbf{Registro de auditoría y evento pendiente.} Conservan los
    cambios del servicio con usuario, fecha y operación, y los eventos que
    aún deben publicarse en RabbitMQ.
\end{itemize}

En el monolito, estas entidades se reparten entre los esquemas
\texttt{usuarios}, \texttt{academico} y \texttt{organizacion} de una base
compartida. En el diseño pasan a una base exclusiva del servicio. La tabla de
auditoría que hoy está en \texttt{usuarios} deja de ser central: cada servicio
registra sus cambios. Las reservas y los laboratorios solo se referencian
mediante identificadores externos. La base por servicio es el patrón de
despliegue de datos más extendido en la práctica, a cambio de resolver la
consistencia entre servicios fuera de la base \cite{laigner2021datamanagement}.

\subsubsection{Componentes internos}

\begin{description}
    \item[API REST.] Recibe las solicitudes del Gateway y las consultas
    internas de otros servicios. Valida la entrada y la dirige al componente
    correspondiente.
    \item[Autenticación y emisión de tokens.] Verifica las credenciales,
    emite un JWT firmado con el identificador del usuario, su rol y sus
    permisos, y registra la sesión activa.
    \item[Gestión de usuarios, roles y permisos.] Crea, activa o desactiva
    cuentas y asigna roles y permisos.
    \item[Solicitudes de acceso.] Registra las peticiones de nuevos usuarios
    y gestiona su aprobación o rechazo.
    \item[Gestión académica.] Administra carreras, asignaturas, horarios,
    periodos y la organización institucional.
    \item[Repositorios de persistencia.] Leen y escriben los datos de la
    base PostgreSQL del servicio.
    \item[Registro de auditoría.] Guarda cada cambio con usuario, fecha y
    operación, junto con el evento pendiente, en la misma transacción.
    \item[Publicador de eventos.] Envía los eventos pendientes a RabbitMQ,
    registra su publicación al recibir la confirmación y reintenta los que
    no se confirmaron.
\end{description}

\subsubsection{Mecanismos de comunicación}

\paragraph{REST sobre HTTPS, sincrónico.}
Los clientes web y móvil ingresan por el API Gateway. El inicio de sesión
necesita una respuesta inmediata con el token o el rechazo, por lo que usa
REST. El Gateway valida la firma y la vigencia del JWT en cada solicitud y
envía a los servicios el usuario, el rol y los permisos; por eso Reservas y
Agenda solo consulta a este servicio para verificar la sesión antes de
confirmar, cancelar o bloquear. Este esquema sigue los
patrones de token de acceso e identidad y de autorización en el borde, en los
que el Gateway concentra la autenticación de las solicitudes
\cite{waseem2022decision,haindl2024interservice}. Laboratorios y Recursos
sí lo consulta al asignar un encargado, porque necesita el resultado antes de
guardar la asignación.

\begin{description}
    \item[Autenticación.] \texttt{POST /auth/login},
    \texttt{POST /auth/logout} y \texttt{GET /auth/mis-roles}.
    \item[Usuarios y accesos.] \texttt{GET /usuarios},
    \texttt{POST /usuarios}, \texttt{PATCH /usuarios/\{id\}/estado},
    \texttt{POST /solicitudes} y \texttt{PATCH /solicitudes/\{id\}}.
    \item[Roles y permisos.] \texttt{GET /roles}, \texttt{POST /roles} y
    \texttt{PUT /roles/\{id\}/permisos}.
    \item[Académico.] \texttt{GET} y \texttt{POST} sobre
    \texttt{/carreras}, \texttt{/asignaturas}, \texttt{/horarios} y
    \texttt{/periodos}.
    \item[Consulta interna.] \texttt{GET /usuarios/\{id\}/roles}. Devuelve
    el estado y los roles de un usuario para validar a un encargado de
    laboratorio.
        \item[Consulta interna.] \texttt{GET /usuarios/\{id\}/roles}. Devuelve
    el estado y los roles de un usuario para validar a un encargado de
    laboratorio.
    \item[Verificación de sesión.] \texttt{GET /auth/sesiones/vigente}.
    Confirma que la sesión no se ha revocado antes de confirmar, cancelar
    o bloquear (RNF-03).
\end{description}

\paragraph{RabbitMQ sobre AMQP, asincrónico.}
El servicio publica los eventos \texttt{UsuarioDesactivado},
\texttt{RolUsuarioModificado}, \texttt{AcademicoActualizado},
\texttt{PeriodoAcademicoCerrado} y \texttt{AccionAuditada}. Cada evento incluye su identificador, el
identificador del recurso afectado, el usuario responsable y la fecha.
Información y Avisos los consume para notificar a los afectados y para la
consulta central de auditoría. El servicio no consume eventos.

\subsubsection{Fallos y consistencia}

Si el servicio no responde, las consultas con un token vigente continúan,
pero las confirmaciones, cancelaciones y bloqueos no se ejecutan, porque
no puede verificarse que la sesión siga activa (RNF-03). Fallan los nuevos
inicios de sesión y la asignación de encargados, que queda sin guardarse. Fallan los nuevos
inicios de sesión y la asignación de encargados, que queda sin guardarse.
Un cambio de rol se aplica en el siguiente token emitido; el vencimiento
corto del JWT limita el tiempo en que un permiso retirado sigue vigente.
El registro local de eventos permite reintentar la publicación si
RabbitMQ se interrumpe, y los consumidores descartan los identificadores
ya procesados.